% 論文節の日本語訳（読み物 PDF）用 preamble
% lualatex + jlreq
\documentclass[12pt,paper=a4]{jlreq}
\usepackage{luatexja-fontspec}
\usepackage{amsmath,amssymb,bm}
\usepackage{graphicx}
\usepackage{hyperref}
\usepackage{enumitem}
\usepackage{booktabs}
\usepackage{microtype}
\usepackage{xcolor}
\usepackage{caption}
\usepackage{tcolorbox}
\tcbuselibrary{breakable,skins}

\hypersetup{
  colorlinks=true,
  linkcolor=blue!60!black,
  urlcolor=blue!60!black,
  citecolor=blue!60!black,
}

\graphicspath{{images/}{./}}

\setlength{\parindent}{1em}
\setlength{\parskip}{0.55em}
\setlist[itemize]{leftmargin=1.6em,itemsep=0.25em,topsep=0.35em}
\setlist[enumerate]{leftmargin=1.8em,itemsep=0.25em,topsep=0.35em}

\captionsetup{
  font=small,
  labelfont=bf,
  skip=0.4em,
  justification=raggedright,
  singlelinecheck=false,
}

\newcommand{\vx}{\bm{x}}
\newcommand{\vu}{\bm{u}}
\newcommand{\veps}{\bm{\epsilon}}
\newcommand{\vw}{\bm{w}}

\DeclareMathOperator*{\argmax}{arg\,max}
\DeclareMathOperator*{\argmin}{arg\,min}
\DeclareMathOperator{\Tr}{Tr}
\DeclareMathOperator{\diag}{diag}

% 命題・定理ボックス（ページ跨ぎ時も切れ目を抑える）
\newtcolorbox{propositionbox}[1][]{%
  enhanced,
  breakable,
  colback=blue!4!white,
  colframe=blue!45!black,
  coltitle=white,
  fonttitle=\bfseries\small,
  title={#1},
  attach boxed title to top left={yshift=-2mm, xshift=3mm},
  boxed title style={
    colback=blue!45!black,
    colframe=blue!45!black,
    boxrule=0pt,
    arc=1pt,
  },
  boxrule=0.5pt,
  arc=1.5pt,
  left=5pt, right=5pt, top=8pt, bottom=5pt,
  before skip=0.9em,
  after skip=0.9em,
  pad at break=3mm,
}

% 図はフロートせず本文位置に置く（分断を防ぐ）
\newcommand{\translationfigure}[2]{%
  \begin{center}
    \includegraphics[width=0.88\linewidth]{#1}\\[0.35em]
    {\small #2\par}
  \end{center}
  \vspace{0.4em}%
}

\begin{document}
{\color{gray}\footnotesize 出典: §5.3\par}
\vspace{0.2em}
\begin{propositionbox}[原題]
\textbf{原題}: Constructing Probability Paths and Velocities Between Distributions
\textbf{日本語タイトル}: 分布間の確率経路と速度の構成
\textbf{出典§}: §5.3 of *The Principles of Diffusion Models*
\end{propositionbox}


\section*{5.3 分布間の確率経路と速度の構成}


フローマッチングの本質は、ソース分布からターゲットへの漸進的変換にある。この変換を方向づけるには、2つの鍵となる要素が必要である：各時刻 $t$ での発展する分布のスナップショットを与える確率経路 $p_t$、および個々の粒子が経路に沿ってどう動くかを記述する速度場 $v_t$ である。これら2つの対象は独立ではなく、粒子ダイナミクスが分布の発展と整合することを保証する連続の式を通じて結び付いている。したがって学習課題は、過程を忠実に駆動する速度場 $v_t$ を見つけることに帰着する。しかし困難は、一般的で複雑な分布に対し真の周辺速度 $v_t$ が未知であり、直接アクセスできない扱いづらい目標が残ることである。
条件付きフローマッチングの核心的着想は、人工的だが扱いやすい過程を構成することで、真の周辺速度の扱いづらさに対処することである。そのために条件付け変数 $z$ を導入し、条件付き速度 $v_t(x_t|z)$ および／または条件付き経路 $p_t(x_t|z)$ を、意図的に単純になるよう設計する。
これらの条件付き対象は閉形式で既知であるため、モデルが回帰できる代理目標として機能する。これは妥当な訓練損失 $\mathcal{L}_{\mathrm{CFM}}$ へ至るが、2つの実務的要件を満たすことが前提である：(i) $p_t(\cdot|z)$ から効率的にサンプリングできること、(ii) 対応する速度 $v_t(\cdot|z)$ が閉形式表現をもつこと。
よく振る舞う条件付き過程をどう設計すべきか？着想を得るため、完全に理解されている唯一の場合——ガウスからガウスへの橋（§5.3.1）——に目を向ける。この例は2つの自然な設計戦略を際立たせる：各時刻 $t$ でガウス確率経路を採用するか、アフィン速度場を規定するかであり、いずれも解析的に扱いやすい。
この洞察に導かれ、2つの相補的な構成視点を用いて一般の端点分布へ拡張する（§B.1.2も参照）。これらは輸送の古典的なオイラー記述とラグランジュ記述に対応する：

\begin{itemize}
\item \textbf{条件付き確率経路優先（オイラー視点）。} 条件付き密度経路 $p_t(\cdot|z)$ から始め、連続の式を満たす条件付き速度場 $v_t(\cdot|z)$ を求める。密度経路だけでは一般に一意な速度を定めないため、この構成は特定の整合する速度場を選ぶことも要する。
\item \textbf{条件付きフロー優先（ラグランジュ視点）。} 代わりに、条件付きフロー写像 $\Psi_{0\to t}(\cdot|z)$ を通じて個々の粒子がどう動くかを指定する。誘導密度は押し出しにより得られ、フロー写像が可逆なとき対応するオイラー速度場は
\end{itemize}

\[
v_t=\partial_t\Psi_{0\to t}\circ\Psi_{0\to t}^{-1}
\]

  である。

§5.3.2では第1のアプローチを詳述し、§5.2.1で議論した拡散モデル構成との密接な類似を示す。§5.3.3では第2を提示する。合わせてこれらの視点は、$p_t(x_t|z)$ と $v_t(x_t|z)$ を定義する実務的枠組みを与え、シミュレーション不要の学習と任意のソース・ターゲット分布間のフロー構成を可能にする。

\subsection*{5.3.1 鍵となる特別な場合：ガウスからガウスへの橋における周辺 $p_t(x_t)$ と速度 $v_t(x_t)$}


まずガウス端点の場合から始める。ここでは周辺密度 $p_t(x_t)$ と速度場 $v_t(x_t)$ を解析的に計算できる。これは一般の条件付き密度 $p_t(x_t|z_t)$ と速度場 $v_t(x_t|z_t)$ の構成の雛形となる。
ソースとターゲット分布 $p_{\mathrm{src}}$ と $p_{\mathrm{tgt}}$ がともにガウスのとき、速度場 $v_t(\cdot)$ は閉形式表現をもつ。補間された周辺密度経路

\[
p_t(x_t)=\mathcal{N}\bigl(x_t;\mu(t),\sigma^2(t)I\bigr) \quad (5.3.1)
\]


を考える。時間変化する平均 $\mu(t)$ と分散 $\sigma^2(t)>0$ をもつ。2つの端点は

\[
p_{\mathrm{src}}=p_0=\mathcal{N}\bigl(x;\mu(0),\sigma^2(0)I\bigr),\qquad
p_{\mathrm{tgt}}=p_1=\mathcal{N}\bigl(x;\mu(1),\sigma^2(1)I\bigr)
\]


で与えられ、経路 $\{p_t\}_{t\in[0,1]}$ がこれらの分布を結ぶ。
与えられた経路 $\{p_t\}_{t\in[0,1]}$ に対し、$x\sim p_0$ ならば $\Psi_{0\to t}(x)\sim p_t$ となるODEフロー $\Psi_{0\to t}(x)$ を誘導する速度場は確かに多数ある。このガウス経路に対する特に単純な実現は次で与えられる\footnote{(Lipman et al., 2022) では著者らは $\Psi_{0\to t}(x)=\mu(t)+\sigma(t)x$ を考え、境界条件を保証するために $\mu(t)$ と $\sigma(t)$ への制約を要する。我々はそのような制約を避ける同値な正規化定式化を採用する。}：

\[
\Psi_{0\to t}(x):=\mu(t)+\sigma(t)\Biggl(\frac{x-\mu(0)}{\sigma(0)}\Biggr). \quad (5.3.2)
\]


定義されたガウス経路 $p_t$（すべての $t$ でガウス）に対し、式 (5.3.2) のODEフローを誘導する速度場 $v_t(\cdot)$ は、一意かつ解析的に次のように特徴づけられる（Lipman et al., 2022）：


\begin{propositionbox}[命題 5.3.1（ガウス密度経路に対する閉形式速度場）]
$p_t$ を式 (5.3.1) のガウス経路とする。するとODEフロー式 (5.3.2) を生成する速度場 $v_t(\cdot)$ は、定義された $\Psi_{0\to t}$ に対して一意であり、閉形式

\[
v_t(x)=\frac{\sigma'(t)}{\sigma(t)}\bigl(x-\mu(t)\bigr)+\mu'(t)
\]

をもつ。
\end{propositionbox}


\textbf{証明の概要。} 初期条件 $y$ をもつODE

\[
\frac{d}{dt}\Psi_{0\to t}(y)=v_t\bigl(\Psi_{0\to t}(y)\bigr)
\]


を考える。$\Psi_{0\to t}$ は可逆（$\sigma(t)>0$）なので、$x=\Psi_{0\to t}(y)$ および $y=\Psi_{0\to t}^{-1}(x)=\Psi_{t\to 0}(x)$ と置いて

\[
\Psi'_{0\to t}\bigl(\Psi_{0\to t}^{-1}(x)\bigr)=v_t(x)
\]


を得る。式 (5.3.2) を $t$ について微分すると

\[
\Psi'_{0\to t}(x)=\mu'(t)+\sigma'(t)\Biggl(\frac{x-\mu(0)}{\sigma(0)}\Biggr).
\]


$y=\Psi_{0\to t}^{-1}(x)$ を解くと

\[
y=\mu(0)+\sigma(0)\Biggl(\frac{x-\mu(t)}{\sigma(t)}\Biggr).
\]


これを $\Psi'_{0\to t}(x)$ に代入すると、主張どおり

\[
v_t(x)=\frac{\sigma'(t)}{\sigma(t)}\bigl(x-\mu(t)\bigr)+\mu'(t)
\]


が得られる。■

固定されたフロー写像 $\Psi_{0\to t}$（フロー優先視点）に対し、速度は

\[
v_t=\partial_t\Psi_{0\to t}\circ\Psi_{0\to t}^{-1}
\]


により一意に定まることに注意する。この構成の下で、対 $(p_t,v_t)$ は自動的に連続の式を満たす。対照的に、$\Psi_{0\to t}$ を固定せずに与えられた密度経路 $t\mapsto p_t$（確率経路優先視点）に対しては、速度場は一意ではない。
この区別がまさに、フロー優先と確率経路優先の視点の差を特徴づける。
この閉形式の特徴づけは、潜在変数 $z$ への条件付けの下でも有効である。以下ではこの洞察を拡張し、一般の周辺設定に対する条件付きガウス経路 $p_t(\cdot|z)$ を構成し、対応する条件付き速度場 $v_t(\cdot|z)$ を導く。

\subsection*{5.3.2 $v_t(\cdot|z)$ と $p_t(\cdot|z)$ の条件付き確率経路優先構成}


\translationfigure{fig-5-9.jpg}{図5.9: 条件付け確率経路の2つのよくある型の模式図。(i) 両側：一般端点分布をもつ $x_0\sim p_{\mathrm{src}}$ と $x_1\sim p_{\mathrm{tgt}}$ に条件付ける；(ii) 片側：$x_0\sim p_{\mathrm{src}}$ または $x_1\sim p_{\mathrm{tgt}}$ のいずれか一端に条件付ける。}


*図5.9: 条件付け確率経路の2つのよくある型の模式図。(i) 両側：一般端点分布をもつ $x_0\sim p_{\mathrm{src}}$ と $x_1\sim p_{\mathrm{tgt}}$ に条件付ける；(ii) 片側：$x_0\sim p_{\mathrm{src}}$ または $x_1\sim p_{\mathrm{tgt}}$ のいずれか一端に条件付ける。*
出典: Created by the authors.

まず条件付き密度経路 $p_t(\cdot|z)$ を構成し、次いで対応する条件付き速度場 $v_t(\cdot|z)$ を（命題5.3.1経由で）導くことを目指す。条件付けは $\pi(z)$ に関して行う。$z$ の選び方に応じて、2つの自然なシナリオがある：(i) $z=(x_0,x_1)$ の両側条件付け、または (ii) $z=x_0$ または $x_1$ の片側条件付け。いずれの場合も、構成は境界分布

\[
p_{\mathrm{src}}(x_0)=\int p_0(x_0|z)\,\pi(z)\,dz,\qquad
p_{\mathrm{tgt}}(x_1)=\int p_1(x_1|z)\,\pi(z)\,dz
\]


と整合しなければならない。具体的構成を指定すればこれらの制約の検証は容易なので、ここでは検証ステップを強調しない。

\subsubsection*{I. 両側 $z=(x_0,x_1)$ ——「ビーム様」経路}


**$\pi(z)$ の選択。** $\mathbb{R}^D$ 上の一般分布 $p_{\mathrm{src}}$ と $p_{\mathrm{tgt}}$ を考える。$z=(x_0,x_1)$ とし、$x_0\sim p_{\mathrm{src}}$、$x_1\sim p_{\mathrm{tgt}}$ を独立に取る。すなわち

\[
\pi(z)=p_{\mathrm{src}}(x_0)\,p_{\mathrm{tgt}}(x_1).
\]


**条件付き経路 $p_t(\cdot|z)$ の選択。** 固定された小さな分散 $\sigma>0$ をもつ線形補間により条件付き経路を定義する\footnote{代わりに固定条件付き分散 $\sigma^2 I$ を加えてもよい。分散が時間一定なので条件付き速度はなお $a_t'x_0+b_t'x_1$ のままであるが、端点周辺は $p_{\mathrm{src}}$ と $p_{\mathrm{tgt}}$ のガウス平滑化となり、$\sigma\to 0$ で望ましい端点に近づく。}：

\[
p_t\bigl(x_t|z=(x_0,x_1)\bigr)=\mathcal{N}(x_t;a_t x_0+b_t x_1,\sigma^2 I),
\]


ここで $a_t$ と $b_t$ は $a_0=1,\ b_0=0$ および $a_1=0,\ b_1=1$ を満たす時間依存関数である。(Lipman et al., 2022; Liu, 2022) が提案する選択は $a_t=1-t$，$b_t=t$ である。決定論的場合 $\sigma=0$ では

\[
p_t(x_t|z)=\delta\bigl(x_t-[a_t x_0+b_t x_1]\bigr)
\]


となり、$x_0$ から $x_1$ への決定論的補間経路を記述する。


**導出される条件付き速度 $v_t(\cdot|z)$。** 命題5.3.1により、条件付き速度は

\[
v_t(x|z)=a_t' x_0+b_t' x_1
\]


である。

\textbf{CFM損失。} $\sigma=0$ で $x_t=a_t x_0+b_t x_1$ のとき、CFM損失は

\[
\mathcal{L}_{\mathrm{CFM}}=\mathbb{E}_{t,\,x_0\sim p_{\mathrm{src}},\,x_1\sim p_{\mathrm{tgt}}}\bigl\|v_\phi(x_t,t)-(a_t'x_0+b_t'x_1)\bigr\|^2
\]


に簡約される。
式 (5.2.9) と (5.2.11) より、最適速度場は

\[
v^*(x_t,t)=\mathbb{E}\bigl[x_t'|x_t\bigr]=\mathbb{E}\bigl[a_t'x_0+b_t'x_1\bigm|x_t\bigr]
\]


である。ここで期待値は $p(x_0,x_1|x_t)$——観測された補間 $x_t=a_t x_0+b_t x_1$ を時刻 $t$ で生じ得たソース–ターゲット対 $(x_0,x_1)$ 上の条件付き分布——上で取る。

\subsubsection*{II. 片側 $z=x_0$ または $x_1$ ——「スポットライト様」経路}


片側設定における条件付き確率経路優先構成を説明する。標準的な生成設定 $p_{\mathrm{src}}=\mathcal{N}(0,I)$、$p_{\mathrm{tgt}}=p_{\mathrm{data}}$ を考える。決定的に重要なのは、このガウスソースは追加仮定ではなく、下記で定義する条件付き経路の直接の帰結であることである。任意端点のより一般的な扱いは§5.3.3で与える。

**$\pi(z)$ の選択。** $z=x_1$ とし $\pi(z)=p_{\mathrm{data}}(x_1)$ を取る（$z=x_0\sim p_{\mathrm{prior}}$ の場合も同様）。

**条件付き経路 $p_t(\cdot|z)$ の選択。** 固定された $x_1\sim p_{\mathrm{data}}$ に対し

\[
p_t(x_t|z=x_1)=\mathcal{N}\bigl(x_t;b_t x_1,a_t^2 I\bigr)
\]


と定義する。$a_0=1,\ b_0=0,\ a_1=0,\ b_1=1$（通常は極限として解釈する）。境界では

\[
p_0(\cdot|z=x_1)=\mathcal{N}(\cdot;0,I),\qquad p_1(\cdot|z=x_1)=\delta(\cdot-x_1).
\]


$x_1$ 上で周辺化すると $\{p_t\}_{t\in[0,1]}$ が得られ、$p_0=\mathcal{N}(0,I)$（$p_{\mathrm{data}}$ に独立）および $p_1=p_{\mathrm{data}}$ となる。

**導出される条件付き速度 $v_t(\cdot|z)$。** $t\in(0,1)$ で $a_t>0$ のとき、条件付きガウス経路に命題5.3.1を適用すると

\[
v_t(x|x_1)=b_t'x_1+\frac{a_t'}{a_t}\bigl(x-b_t x_1\bigr)
\]


が得られる。

\textbf{片側CFM目的。} $t\sim\mathcal{U}(0,1)$（または任意の固定サンプリング分布）および $x_1\sim p_{\mathrm{data}}$ の下で、CFM損失は

\[
\mathcal{L}_{\mathrm{CFM}}=\mathbb{E}_{t,x_1}\,\mathbb{E}_{x_t\sim p_t(\cdot|x_1)}\Biggl\|v_\phi(x_t,t)-\Biggl(b_t'x_1+\frac{a_t'}{a_t}(x_t-b_t x_1)\Biggr)\Biggr\|_2^2 \quad (5.3.3)
\]


となる。MSE最適性により、一意な最小化子は周辺速度場

\[
v^*(x,t)=\mathbb{E}\bigl[v_t(x|x_1)\bigm|x_t=x\bigr]=\mathbb{E}\bigl[a_t'x_0+b_t'x_1\bigm|x_t=x\bigr]
\]


である。

\textbf{両側目標との同値性。} 対付けサンプル $(x_0,x_1)$ で $x_t=a_t x_0+b_t x_1$ のとき、

\[
v_t(x_t|x_1)=b_t'x_1+\frac{a_t'}{a_t}(x_t-b_t x_1)=a_t'x_0+b_t'x_1.
\]


したがって片側損失は、$x_t$ が与えられたときの両側目標の条件付き期待値へ回帰する：

\[
v^*(x,t)=\mathbb{E}\bigl[a_t'x_0+b_t'x_1\bigm|x_t=x\bigr],
\]


ゆえに片側と両側のCFM目的は同じ最小化子を共有する。

**ガウスFM $=$ 拡散モデル。** FMの慣例では $t=0$ がソース／事前、$t=1$ がターゲット／データを表す：

\[
p_{\mathrm{src}}=p_{\mathrm{prior}},\qquad p_{\mathrm{tgt}}=p_{\mathrm{data}}.
\]


対照的に、拡散モデルは典型的に時間をデータからノイズへ索引付けする（すなわち $t=0$ がデータ、$t=1$ が事前）。ここでは主議論でFM慣例の索引を採用する。ただし比較のため、表5.2のVEおよびVPの項目は通常の拡散時間索引（$t=0$ がデータ、$t=1$ が事前／ノイズ）で要約し、$a_t:=\sigma_t$、$b_t:=\alpha_t$ とする。さらに $p_{\mathrm{src}}=\mathcal{N}(0,I)$ なら、固定条件 $x_1\sim p_{\mathrm{data}}$ に対し条件付き経路 $p_t(\cdot|x_1)$ は自然にガウスに選ばれ、ターゲット分布 $p_{\mathrm{tgt}}$ 自体はガウスである必要はない。このガウス条件付き経路設定をガウスFMと呼ぶ。拡散確率経路は、そのようなガウス経路の重要な部分クラスをなす。しかし一般のガウススケジュールは、妥当な前方拡散SDEに対応する必要はない；スケジュールはさらに補題4.5.1のSDE整合条件を満たさねばならない。
$a_t=1-t$、$b_t=t$ を選ぶと（同値に、拡散モデルでの再ラベル $a_t:=\sigma_t$、$b_t:=\alpha_t$ の下で $\alpha_t=t$、$\sigma_t=1-t$）、おなじみのFM/RFスケジュール（Lipman et al., 2022; Liu, 2022）が回復される。

\translationfigure{table-5-2.jpg}{表5.2: 異なる補間子の要約。$x_t=a_t x_0+b_t x_1$ と書き、$x_0\sim p_{\mathrm{src}}=p_{\mathrm{prior}}$、$x_1\sim p_{\mathrm{tgt}}=p_{\mathrm{data}}$。FM/RFおよび三角法列はFM慣例（$t=0$ ソース、$t=1$ ターゲット）を用い、VE/VP列は通常の拡散時間慣例（$t=0$ データ、$t=1$ 事前／ノイズ）を保持し、$a_t:=\sigma_t$、$b_t:=\alpha_t$。有限終端ノイズ水準のVEでは、厳密な終端周辺は $p_{\mathrm{data}}*\mathcal{N}(0,a_1^2 I)$ であり、表に示すガウス事前は終端ノイズがデータの尺度を支配するときの通常の近似である。}


*表5.2: 異なる補間子の要約。$x_t=a_t x_0+b_t x_1$ と書き、$x_0\sim p_{\mathrm{src}}=p_{\mathrm{prior}}$、$x_1\sim p_{\mathrm{tgt}}=p_{\mathrm{data}}$。FM/RFおよび三角法列はFM慣例（$t=0$ ソース、$t=1$ ターゲット）を用い、VE/VP列は通常の拡散時間慣例（$t=0$ データ、$t=1$ 事前／ノイズ）を保持し、$a_t:=\sigma_t$、$b_t:=\alpha_t$。*

ガウスFM設定では、ビーム様・スポットライト様の両方の条件付き経路が、標準的な拡散損失に類似した学習目的を生じる。第6章で詳述するように、ガウスFMは実際、線形スケジュール $a_t=1-t$、$b_t=t$ の下で速度を予測するよう訓練された拡散モデルと同値に解釈できる。この視点は、フローマッチングと拡散が根本的に異なるのではなく、相互に変換可能な2つの同値な定式化であることを際立たせる。ガウスFM目的は実務上特に魅力的である：その損失関数（$\mathbb{E}_{t,x_t}\|v_\phi(x_t,t)-(x_1-x_0)\|_2^2$）は単純であり、大規模で競争力のある性能を達成することが示されている（Esser et al., 2024）。

\subsection*{5.3.3 $v_t(\cdot|z)$ と $p_t(\cdot|z)$ の条件付きフロー優先構成}


端点 $p_{\mathrm{src}}$（$t=0$）と $p_{\mathrm{tgt}}$（$t=1$）が任意である一般の場合を扱う。目標は、軌道空間で直接、$p_{\mathrm{src}}$ から $p_{\mathrm{tgt}}$ へサンプルを輸送し、回帰目標として使える閉形式の $v_t(x_t|z)$ を生じる条件付きフローを設計することである。

\textbf{動機。} まず条件付き密度経路を設計する代わりに、軌道に沿ってサンプルを動かす条件付きフロー写像 $\Psi_{0\to t}(\cdot;z)$ を直接指定してもよい。これには2つの実務的利点がある：(i) 軌道に沿った時間微分を通じて訓練用の回帰目標を直ちに与える；(ii) 幾何構造をもつ空間（リーマン多様体、リー群、または制約付き部分多様体）では、幾何から直接条件付きフロー写像 $\Psi_{0\to t}$ を構成するのがしばしば自然である（例：測地線、指数写像、または擬距離）（Lipman et al., 2024）。これにより訓練用の解析的・シミュレーション不要の目標速度が得られる。

\textbf{条件付きアフィンフロー（命題5.3.1へのリンク）。} 条件付け変数 $z\sim\pi$（例：片側「スポットライト」訓練での $z=x_1\sim p_{\mathrm{tgt}}$）を固定し、$x_0\sim p_{\mathrm{src}}$ を時間変化する条件付きアフィンフロー

\[
\Psi_{0\to t}(x_0;z):=\mu_t(z)+A_t(z)x_0,\qquad t\in[0,1]
\]


により押し出す。ここで $\mu_t(z)\in\mathbb{R}^D$、$A_t(z)\in\mathbb{R}^{D\times D}$ は $t\in(0,1)$ で可逆である。境界 $A_0(z)=I$、$\mu_0(z)=0$ が $t=0$ で $p_{\mathrm{src}}$ を回復する。$t\to 1$ の境界は極限として解釈するのが標準である（終端写像はより低次元の集合または点に質量を集中させ得る）\footnote{$A_1(z)$ が特異（例：$0$）であることを許すのは、$(0,1)$ 上の可逆性と両立し、$t=1$ で経路が規定端点へ収縮する。}。


**誘導される条件付き経路 $p_t(\cdot|z)$。** 構成は

\[
p_t(\cdot|z)=\Psi_{0\to t}(\cdot;z)\# p_{\mathrm{src}},\qquad
p_t(\cdot)=\int p_t(\cdot|z)\,\pi(z)\,dz
\]


を定義する。$\mathcal{L}_{\mathrm{CFM}}$ で最終的に重要なのは、それからどうサンプリングするかである：まず $z\sim\pi$ を引き、次いで $x_0\sim p_{\mathrm{src}}$ を引き、最後に

\[
x_t=\mu_t(z)+A_t(z)x_0
\]


と置く。
$\Psi_{0\to t}$ が $x_0$ に関してアフィンなとき、$p_t(\cdot|z)$ がガウスであるのは $p_{\mathrm{src}}$ がガウスのとき、かつそのときに限ることに注意する。特に任意の（非ガウス）$p_{\mathrm{src}}$ に対し、アフィンフローは一般に非ガウスな $p_t(\cdot|z)$ を生じる。

**導出される条件付き速度 $v_t(\cdot|z)$。** 条件付き速度 $v_t(\cdot|z)$ は、条件付きフロー写像 $\Psi_{0\to t}$ を $t$ 微分することで得られる。命題5.3.1の導出に従い、初期条件 $y$ をもつフロー写像 $\Psi_{0\to t}(y;z)$ で定義される条件付きODEを考え、対応する条件付き速度場 $v_t(\cdot|z)$ を同定する：

\[
\frac{d}{dt}\Psi_{0\to t}(y;z)=v_t\bigl(\Psi_{0\to t}(y;z)\bigm|z\bigr).
\]


$\Psi_{0\to t}(\cdot;z)$ は $t\in(0,1)$ で可逆なので、現在状態 $x=\Psi_{0\to t}(y;z)$ に対応する初期点 $y$ を

\[
y=\Psi_{0\to t}^{-1}(x;z)=\Psi_{t\to 0}(x;z)
\]


と表せる。対応するオイラー速度場はしたがって

\[
v_t(x|z):=\frac{\partial}{\partial t}\Psi_{0\to t}(y;z)\Bigm|_{y=\Psi_{0\to t}^{-1}(x;z)}
\]


である。すなわち、まず初期点 $y$ を固定したまま前方フロー写像を $t$ について微分し、次いで時刻 $t$ で $x$ に到達する初期点を代入する。
$x_t=\mu_t(z)+A_t(z)x_0$ かつ $A_t(z)$ が $(0,1)$ 上で可逆なので、$x_0=A_t(z)^{-1}\bigl(x-\mu_t(z)\bigr)$ となり、

\[
v_t(x|z)=\mu_t'(z)+A_t'(z)A_t(z)^{-1}\bigl(x-\mu_t(z)\bigr)
\]


が得られる。

**片側条件付け（$z=x_1$）。** $\mu_t(z)=b_t z$、$A_t(z)=a_t I$ を選び、$a_0=1$、$a_1=0$、$b_0=0$、$b_1=1$（$t\in(0,1)$ で $a_t>0$）とすると

\[
x_t=a_t x_0+b_t x_1,\qquad
v_t(x|x_1)=b_t'x_1+\frac{a_t'}{a_t}\bigl(x-b_t x_1\bigr)
\]


が得られる。対付けサンプル $(x_0,x_1)$（$x_t=a_t x_0+b_t x_1$）では、通常のCFM目標

\[
v_t(x_t|x_1)=a_t'x_0+b_t'x_1
\]


に簡約される。

**両側条件付け（$z=(x_0,x_1)$）。** 同じ雛形で $\mu_t(x_0,x_1)=b_t x_1$、$A_t(x_0,x_1)=a_t I$ とすると、条件付き経路は決定論的になる：

\[
x_t=a_t x_0+b_t x_1,\qquad p_t(\cdot|x_0,x_1)=\delta\bigl(\cdot-(a_t x_0+b_t x_1)\bigr),
\]


条件付き速度は

\[
v_t(x_t|x_0,x_1)=a_t'x_0+b_t'x_1,
\]


すなわち標準的な両側CFM目標である。

\textbf{無条件ガウス経路を特別な場合として。} $\mu_t$ が $z$ に独立（$\mu(t)$ と表記）で $A_t=\frac{\sigma(t)}{\sigma(0)}I$ なら、

\[
\Psi_{0\to t}(x_0)=\mu(t)+\sigma(t)\frac{x_0-\mu(0)}{\sigma(0)},\qquad
v_t(x)=\mu'(t)+\frac{\sigma'(t)}{\sigma(t)}\bigl(x-\mu(t)\bigr)
\]


となり、命題5.3.1のガウス密度経路と閉形式速度が回復される。

\subsection*{5.3.4 確率経路優先構成とフロー優先構成の比較}


2つの構成は、輸送の相補的なオイラー記述とラグランジュ記述に対応する。確率経路優先（オイラー）視点は、まず条件付き密度経路 $p_t(\cdot|z)$ を指定し、次いで連続の式を通じて整合する速度場を選ぶ。密度経路だけでは一般に速度を一意に定めないため、1つを選ぶために追加の輸送構造が必要になり得る。
フロー優先（ラグランジュ）視点は代わりに条件付きフロー写像 $\Psi_{0\to t}(\cdot|z)$ を指定し、粒子軌道を直接記述する。対応する密度経路は押し出しに従い、写像が可逆なときそのオイラー速度は

\[
v_t=\partial_t\Psi_{0\to t}\circ\Psi_{0\to t}^{-1}
\]


として一意に回復される。
したがって両構成は異なる出発点から同じ種類の輸送対象を記述するが、規定された密度経路が、別途規定されたフロー写像と同じ速度を選ぶ必要はない。
次の表がこれらの対比を要約する。要点：経路優先は条件付き経路が閉形式速度をもつときに自然；フロー優先は軌道に強い構造的事前があるときに自然である。

\translationfigure{table-5-3.jpg}{表5.3: 条件付き確率経路優先構成と条件付きフロー優先構成の比較。}


*表5.3: 条件付き確率経路優先構成と条件付きフロー優先構成の比較。与えられるもの、速度の得方、閉形式、$v_t(\cdot|z)$ の一意性、実現可能性、$(p_{\mathrm{src}},p_{\mathrm{tgt}})$ の整合、および好ましいシナリオの各軸で対比される。*

\subsection*{5.3.5 条件付き直線性は周辺フローの直線性を含意しない}


正準アフィン補間

\[
x_t=(1-t)x_0+t x_1
\]


は魅力的な幾何的性質をもつ：特定の端点対 $(x_0,x_1)$ に条件付けると、補間は一定速度 $x_1-x_0$ の直線である。この条件付き直線性の概念は、フローマッチングと整流フローで広く用いられる。
しかし、これらの条件付き補間経路の直線性と、生成時に用いられる周辺ODE軌道の直線性を区別することが重要である。カップリング $\pi(x_0,x_1)$ の下で、母集団最適周辺速度は

\[
v^\star(x,t)=\mathbb{E}\bigl[x_1-x_0\bigm|x_t=x\bigr]
\]


である。端点対上の事後は一般に $(x,t)$ とともに変化する。したがって、訓練に用いる各条件付き補間が直線であっても、周辺ODE軌道に沿って遭遇する速度は一定である必要がない。
こうして、正準アフィンスケジュールだけを選んでも、学習された周辺ODEが厳密に直線であることや、1回のオイラーステップが厳密であることは含意されない。整流フローはまさにこの区別に対処する：その整流または Reflow 手続きは端点カップリングを変え、誘導される因果フローが次第に直線フローに近づける（Liu, Gong, et al., 2022）。
次の例が区別を明示する。

\subsubsection*{例：単純な1次元反例}


1次元設定で

\[
x_0\sim\mathcal{N}(0,\sigma_0^2),\qquad x_1\sim\mathcal{N}(0,\sigma_1^2),
\]


独立カップリング、および正準アフィン補間 $x_t=(1-t)x_0+t x_1$ を考える。各条件付き経路 $(1-t)x_0+t x_1$ は時間について直線である。次に誘導される周辺速度場を計算する。
$(x_1-x_0,x_t)$ が結合ガウスなので、条件付き期待値は $x$ に関して線形であり、

\[
v^*(x,t)=\mathbb{E}[x_1-x_0|x_t=x]=\frac{\mathrm{Cov}(x_1-x_0,x_t)}{\mathrm{Var}(x_t)}\,x.
\]


直接計算により

\[
\mathrm{Var}(x_t)=(1-t)^2\sigma_0^2+t^2\sigma_1^2=:V(t),
\]


および

\[
\mathrm{Cov}(x_1-x_0,x_t)=t\sigma_1^2-(1-t)\sigma_0^2=\tfrac{1}{2}V'(t)
\]


が得られる。ゆえに

\[
v^*(x,t)=\frac{V'(t)}{2V(t)}\,x.
\]


周辺ODE

\[
\frac{dz_t}{dt}=v^*(z_t,t)
\]


はしたがって厳密解

\[
z_t=z_0\sqrt{\frac{V(t)}{V(0)}}=z_0\sqrt{(1-t)^2+\Bigl(\frac{\sigma_1}{\sigma_0}\Bigr)^2 t^2}
\]


をもつ。
さて、時間についての直線補間は $z_t$ が $t$ に関してアフィンであること、すなわちある端点 $z_1$ に対し $z_t=(1-t)z_0+t z_1$ であることを要求する。しかし

\[
\sqrt{(1-t)^2+\frac{\sigma_1^2}{\sigma_0^2}t^2}
\]


は、任意の正の $\sigma_0,\sigma_1$ に対して $t$ に関してアフィンではない。したがってODE軌道は一般に時間についての直線補間ではない。この現象を §5.3.5 で $\sigma_0=2$、$\sigma_1=3$ の例により図示する。
例えば $\sigma_0=\sigma_1$ のとき $z_1=z_0$ となり、直線補間は単に定数経路 $z_t=z_0$ となる。対照的に実際のODE軌道は

\[
z_{1/2}=\frac{z_0}{\sqrt{2}}\neq z_0
\]


を満たす。したがって、すべての条件付き補間経路が直線であっても、誘導される周辺ODE軌道はそうである必要がない。■

\translationfigure{fig-5-10.jpg}{図5.10: 正準アフィン補間の下で周辺ODE軌道は直線ではない。$x_0\sim\mathcal{N}(0,\sigma_0^2)$、$x_1\sim\mathcal{N}(0,\sigma_1^2)$ で $\sigma_0=2$、$\sigma_1=3$、独立カップリングを取る。灰色の線はサンプリングされた条件付き経路 $(1-t)x_0+tx_1$ を示し、構成上直線である。実線は厳密な周辺ODE軌道 $z_t=z_0\sqrt{(1-t)^2+(\sigma_1/\sigma_0)^2 t^2}$、破線は $z_0$ から $z_1$ への対応する直線補間を示す。実線と破線の可視のギャップは、すべての条件付き経路が直線であっても周辺ODE軌道が $t$ に関してアフィンでないことを確認する。周辺密度 $p_{\mathrm{src}}$ と $p_{\mathrm{tgt}}$ はそれぞれ左・右マージンに示される。}


*図5.10: 正準アフィン補間の下で周辺ODE軌道は直線ではない。$x_0\sim\mathcal{N}(0,\sigma_0^2)$、$x_1\sim\mathcal{N}(0,\sigma_1^2)$ で $\sigma_0=2$、$\sigma_1=3$、独立カップリングを取る。灰色の線はサンプリングされた条件付き経路 $(1-t)x_0+tx_1$ を示し、構成上直線である。実線は厳密な周辺ODE軌道 $z_t=z_0\sqrt{(1-t)^2+(\sigma_1/\sigma_0)^2 t^2}$、破線は $z_0$ から $z_1$ への対応する直線補間を示す。実線と破線の可視のギャップは、すべての条件付き経路が直線であっても周辺ODE軌道が $t$ に関してアフィンでないことを確認する。周辺密度 $p_{\mathrm{src}}$ と $p_{\mathrm{tgt}}$ はそれぞれ左・右マージンに示される。*
出典: Created by the authors with AI-assisted coding.
\end{document}
